\section{Introduction}
\label{sec:intro}

An attractor is hidden if its basin does not intersect small neighborhoods of the equilibria, so that it cannot be found
by starting trajectories near an equilibrium; otherwise it is self-excited.\cite{LeonovKuznetsov2013,Dudkowski2016} Hidden
attractors were first located in Chua's circuit\cite{Leonov2011PLA} and are now known in systems without equilibria, with
lines of equilibria and with a single stable equilibrium.\cite{Dudkowski2016,Molaie2013,JafariSprott2013,WangChen2012} The
distinction matters in practice because the standard analysis of a system starts at its equilibria: their stability and
the unstable manifolds of the saddles reveal the self-excited regimes, while a hidden regime stays invisible. In control
systems, phase-locked loops and drilling rigs, such regimes coexist with a stable operating point and appear only after a
large perturbation.\cite{Kuznetsov2020JCSSI}

Models trained from data increasingly stand in for equations. Reservoir computers running in closed loop reproduce chaotic
attractors and their Lyapunov spectra\cite{JaegerHaas2004,Pathak2017Chaos,LuHuntOtt2018} and forecast large spatiotemporal
systems;\cite{Pathak2018PRL} next-generation reservoir computers,\cite{Gauthier2021NGRC} recurrent networks trained by
backpropagation,\cite{Vlachas2020Backprop} sparse regression\cite{Brunton2016SINDy} and neural ordinary differential
equations\cite{Chen2018NeuralODE} serve the same purpose. Beyond forecasting, such models are asked for global information:
the coexisting attractors of a multistable system and their basins,\cite{Roy2022Multistability,Liu2025HiddenDynamics} and
even attractors absent from the training data. A reservoir trained on one attractor of the Li--Sprott system, which has no
equilibria, reproduces its unseen coexisting attractors,\cite{Rohm2021,Rohm2022Unseen} and a reservoir trained on the
trajectories of one basin reproduces other basins and chaotic attractors.\cite{Norton2025Unseen}

A trained model is, however, a dynamical system of its own, often of dimension in the hundreds, and its behavior away from
the data is not controlled by training. Reservoirs acquire mirror attractors from the symmetry of their
equations,\cite{HerteuxRath2020} extra saddles that organize their dynamics,\cite{Kabayama2025PseudoLorenz} untrained
attractors and fixed points,\cite{Flynn2021,Hart2024CLE,Norton2025Unseen} and spurious slow modes that pull trajectories to
fixed points;\cite{Oishi2026FastSlow} untrained attractors have been argued to be intrinsic to learning systems with a
bounded state space.\cite{OHagan2025Confabulation} Generalization across basins is impossible in general without prior
structure.\cite{Goring2024OOD} Most of this literature treats spurious behavior as a symptom of imperfect training, and it
studies attractors that are either visible in the data or have no equilibria nearby.

A hidden attractor of a system with equilibria is the hardest case of this kind. Reproducing it requires the model to get
right not only the attractor but also the neighborhoods of all equilibria, including equilibria that the true system does
not have. A model can reproduce the attractor faithfully and still place an unstable equilibrium whose neighborhood leads to
it; the attractor then becomes self-excited in the model, and the model gives the wrong answer for states near its
equilibria. Reservoir computers have been trained on systems with hidden attractors,\cite{Xu2024SelfESN,Mahata2025Extreme}
and one study predicted the basin of the stable equilibrium of the Wang--Chen system next to what is, at its parameters, a
hidden cycle,\cite{Liu2025HiddenDynamics} but to our knowledge no work has checked whether a reproduced attractor stays
hidden with respect to the equilibria of the model; the closest works are those just cited and the index-theoretic
results for threshold-linear and Hopfield networks discussed in Sec.~\ref{sec:degree}.

This paper asks whether trained models keep hidden attractors hidden, and answers it for about $10^4$ models of four
classes trained on six regimes of four systems with hidden attractors, plus one weakly hidden control; the survival curves
rest on $3710$ reproductions of a hidden attractor, and a full-protocol ensemble of echo state networks adds $581$. The answer is mostly
no. Topology fixes the bookkeeping of a model's
equilibria and, in some systems, forces spurious ones on models that copy the true equilibria; training decides how many unstable equilibria the model has; and
their number in the region of the true system largely decides whether hiddenness survives. We make four contributions.
\begin{enumerate}
\item We adapt classical degree theory to trained models (Theorem~\ref{thm:degree}): a bounded model has an odd number of
equilibria, at least $(K-1)/2$ of them unstable, and an index sum different from $1$ certifies that a numerically found set
of equilibria is incomplete, a check we use throughout.
\item A bounded model that reproduces the true equilibria with their indices must have spurious equilibria of total index
$1-s$, where $s$ is the index sum of the true equilibria (Proposition~\ref{prop:forced}). For Chua's circuit in regime~I,
$s=-1$, and no bounded model in our ensembles reproduces all three equilibria; for the equilibrium-free system of R\"ohm
\textit{et al.},\cite{Rohm2021} $s=0$, every model has a spurious equilibrium, and the reproduced training attractor becomes
self-excited.
\item An attractor of the model is excited from an equilibrium exactly when the local unstable manifold of the equilibrium
reaches the closure of its basin (Proposition~\ref{prop:manifold}), which reduces the excitation test to the unstable
subspace.
\item The probability that a reproduced hidden attractor stays hidden falls with the number of unstable equilibria of the
model inside the region of the true system, similarly for echo state networks, one- and two-layer perceptrons and residual
networks; model class, reservoir size and depth act mainly through this number. Data from the basin of the visible
attractor raise survival. A penalty on unstable equilibria during training, which removes
them at fixed data, raises survival from $0.27$ to $0.994$ at an unchanged one-step error. Generic models do not reproduce a
hidden attractor that the data do not touch; a model whose function class contains the Euler step of the true vector
field sometimes does.
\end{enumerate}

Section~\ref{sec:setting} defines the models, the notion of hiddenness for a model and the test systems.
Section~\ref{sec:theory} states the theory, and Sec.~\ref{sec:methods} describes the ensembles and the labeling protocol.
Section~\ref{sec:results} tests the theory, measures the survival law and the effect of the data, and
Sec.~\ref{sec:discussion} discusses what this means for the use of trained models.
